\documentclass[a4paper,13pt]{article}

\usepackage[UTF8,fontset=none,linespread=1.0]{ctex}
% 仿宋只有常规字重，加粗用伪粗体补齐
\setCJKmainfont{FandolFang-Regular.otf}[AutoFakeBold=3.5]
\setCJKsansfont{FandolFang-Regular.otf}[AutoFakeBold=3.5]
\setCJKmonofont{FandolFang-Regular.otf}[AutoFakeBold=3.5]
\usepackage{geometry}
\usepackage{titlesec}
\usepackage{enumitem}
\usepackage[hidelinks]{hyperref}
\usepackage{fontawesome}
\usepackage{xcolor}
\usepackage{tabularx}

\geometry{left=0.78in,top=0.44in,right=0.78in,bottom=0.5in}

\newcommand{\namefontsize}{\fontsize{22}{26}\selectfont}
\newcommand{\sectiontitlefontsize}{\fontsize{16}{19}\selectfont}
\newcommand{\entrytitlefontsize}{\fontsize{14}{19}\selectfont}
\newcommand{\largeentrytitlefontsize}{\fontsize{15}{19}\selectfont}
\newcommand{\smallentrytitlefontsize}{\fontsize{12}{16}\selectfont}
\newcommand{\entrydetailsfontsize}{\fontsize{11.5}{14}\selectfont}
\newcommand{\contactfontsize}{\fontsize{11}{14}\selectfont}
\newcommand{\linkfontsize}{\fontsize{11.8}{14.6}\selectfont}
\newcommand{\descfontsize}{\fontsize{11}{14.6}\selectfont}
\newcommand{\itemfontsize}{\fontsize{11}{14.6}\selectfont}

\definecolor{applegray}{RGB}{78,78,81}
\definecolor{textgray}{RGB}{45,45,45}
\definecolor{sectiongray}{RGB}{60,60,65}
\definecolor{appledark}{RGB}{0,0,0}
\definecolor{appleblue}{RGB}{0,0,139}

\pagestyle{empty}
\urlstyle{same}
\raggedbottom
\raggedright
\setlength{\tabcolsep}{0in}
\setlength{\parindent}{0pt}
\setlength{\parskip}{0pt}

\titleformat{\section}
    {\sectiontitlefontsize\bfseries\color{sectiongray}}
    {}{0em}{}
\titlespacing*{\section}{0pt}{14pt}{0pt}
\newcommand{\sectionbodyspace}{\vspace{1.5pt}{\color{applegray}\hrule height 0.8pt}\vspace{4pt}}
\newcommand{\sectiongap}{\vspace{-3pt}}
\newcommand{\entrygap}{\vspace{1pt}}
\newcommand{\projectgap}{\vspace{6pt}}

\newcommand{\resumeSubheading}[4]{
    \vspace{0pt}
    \begin{tabular*}{\textwidth}{l@{\extracolsep{\fill}}r}
        \textbf{\entrytitlefontsize\color{appledark}#1} \textit{\entrydetailsfontsize\color{textgray} · #3#4} & {\entrydetailsfontsize\color{applegray}#2} \\
    \end{tabular*}
    \vspace{-2pt}
}

\newcommand{\resumeItemListStart}{
    \begin{itemize}[
        leftmargin=1.5em,
        labelsep=0.5em,
        label={\color{applegray}\raisebox{0.12ex}{\scriptsize\faAngleRight}},
        itemsep=6pt,
        parsep=0pt,
        topsep=1pt
    ]
}
\newcommand{\resumeItemListEnd}{\end{itemize}\vspace{1pt}}
\newcommand{\skillItemListStart}{
    \begin{itemize}[
        leftmargin=1.5em,
        labelsep=0.5em,
        label={\color{applegray}\raisebox{0.12ex}{\scriptsize\faAngleRight}},
        itemsep=3pt,
        parsep=0pt,
        topsep=1pt
    ]
}

\newcommand{\headerLinkItem}[2]{
    \textcolor{textgray}{\linkfontsize#1}~\textcolor{appleblue}{\linkfontsize#2}
}

\newcommand{\resumeItem}[1]{\item\itemfontsize{#1}}

\begin{document}

% --- page 1 (English): same compact engine as resume_EN.tex ---
\renewcommand{\baselinestretch}{1.0}
\titleformat{\section}
    {\sectiontitlefontsize\bfseries\color{sectiongray}}
    {}{0em}
    {}[\vspace{-14pt}{\color{applegray}\rule{\textwidth}{0.8pt}}]
\titlespacing{\section}{0pt}{9pt}{3pt}
\renewcommand{\sectionbodyspace}{\vspace{-3pt}}
\renewcommand{\sectiongap}{\vspace{-7pt}}
\renewcommand{\projectgap}{\vspace{5pt}}
\renewcommand{\resumeItemListStart}{%
    \begin{itemize}[leftmargin=1.5em, labelsep=0.5em,
        label={\color{applegray}\raisebox{0.12ex}{\scriptsize\faAngleRight}},
        itemsep=3pt, parsep=0pt, topsep=1pt]}
\noindent
\begin{minipage}[t]{0.60\textwidth}
\vspace{0pt}
{\namefontsize\bfseries\color{appledark}Longchang Liu}\par\vspace{5pt}
{\entrytitlefontsize\bfseries\color{textgray}Software Engineer}\par\vspace{6pt}
{\contactfontsize\textcolor{textgray}{\faPhone}~\textcolor{textgray}{+86 183-5757-1400}\quad\textcolor{textgray}{\faEnvelope}~\textcolor{textgray}{liulch.cn@gmail.com}}\par
\end{minipage}\hfill
\begin{minipage}[t]{0.36\textwidth}
\vspace{0pt}
\raggedright
\vspace{4pt}
{\headerLinkItem{\faGithub}{\href{https://github.com/leochame}{github.com/leochame}}}\par\vspace{3pt}
{\headerLinkItem{\faLinkedin}{\href{https://www.linkedin.com/in/longchang-liu-770211387/}{linkedin.com/in/longchang-liu}}}\par\vspace{3pt}
{\headerLinkItem{\faGlobe}{\href{https://leochame.github.io}{leochame.github.io}}}\par
\end{minipage}
\par\vspace{2pt}

\section{Education}
\sectionbodyspace
\begin{tabular*}{\textwidth}{l@{\extracolsep{\fill}}r}
    \textbf{\smallentrytitlefontsize\color{appledark}Hangzhou Dianzi University}\hspace{0.8em}{\descfontsize\color{textgray}B.S. in Computer Science and Technology} & {\descfontsize\color{applegray}Aug 2022 - Jun 2026} \\
\end{tabular*}
\sectiongap

\section{Experience}
\sectionbodyspace

\begin{tabular*}{\textwidth}{l@{\extracolsep{\fill}}r}
    \textbf{\smallentrytitlefontsize\color{appledark}Siorigin - Agentic Infra (Kernel-Gen Agent)} & {\entrydetailsfontsize\color{applegray}Jul 2026 - Present} \\
\end{tabular*}
{\descfontsize\color{textgray} Designed the task contract, verification flow and evaluation rubric for a kernel-generation agent that autonomously generates, validates and tunes operators for internal workloads; it reaches the published state of the art on FlashInfer benchmarks.}\vspace{2pt}
\resumeItemListStart
\resumeItem{\textbf{Task contract}: the target operator, constraints and promotion criteria are fixed as the agent's per-round input, so the objective does not drift over a long loop.}
\resumeItem{\textbf{Verification}: after every change, numerical comparison against a PyTorch reference, then scoring on measured latency and achieved bandwidth --- both gates required before the next round.}
\resumeItem{\textbf{Search and context}: candidates, profiling evidence and the promotion decision are recorded each round; failure reasons feed the next round and only useful evidence is kept, bounding context growth. Iteration caps and regression rejection keep the agent from repeating ineffective attempts.}
\resumeItem{\textbf{Delivery}: internal users supply only an operator and a reference implementation, and get back a validated kernel.}
\resumeItemListEnd
\projectgap

\begin{tabular*}{\textwidth}{l@{\extracolsep{\fill}}r}
    \textbf{\smallentrytitlefontsize\color{appledark}Alibaba - Amap (Intern)} & {\entrydetailsfontsize\color{applegray}Oct 2025 - Jan 2026} \\
\end{tabular*}
{\descfontsize\color{textgray} Responsible for productionizing the Order Analysis Assistant and AI Ride-Hailing LLM projects, improving complex task handling and integrating shared AI capabilities.}\vspace{2pt}
\resumeItemListStart
\resumeItem{On the Order Analysis Assistant, designed and launched Deep Research, improving complex task decomposition and handling stability in long-context scenarios.}
\resumeItem{Participated in the cold start of AI Ride-Hailing, building unified and stable AI invocation interfaces to keep the core path reliably integrated.}
\resumeItemListEnd
\projectgap

\begin{tabular*}{\textwidth}{l@{\extracolsep{\fill}}r}
    \textbf{\smallentrytitlefontsize\color{appledark}Tencent - Ads ChatBI System (Intern)} & {\entrydetailsfontsize\color{applegray}Apr 2025 - Sep 2025} \\
\end{tabular*}
{\descfontsize\color{textgray} Natural-language data query system for content business users. Grounded in a metrics platform as its knowledge source, it fills DSL templates step by step through multi-turn clarification and multi-path retrieval with reranker re-ranking and popularity weighting, suppressing hallucination; a bad-case to fine-tuning to ICL hot-fix loop drives continuous improvement, with a state machine providing crash recovery.}\par
\sectiongap

\section{Open Source}
\sectionbodyspace
{\entrydetailsfontsize\color{textgray}\textbf{\color{appledark}LangChain4j} (12k+ stars) contributor --- Java ecosystem AI framework; community issue fixes and capability improvements.}\par
\sectiongap

\section{Technical Skills}
\vspace{0pt}
\skillItemListStart
    \resumeItem{\textbf{AI Infra}: Transformer architecture and KV Cache; agent context management, multi-agent collaboration patterns, harness design and eval systems, with hands-on GPU kernel generation and tuning.}
    \resumeItem{\textbf{CS Fundamentals}: Solid OS and networking fundamentals --- memory management (page tables), system calls, thread scheduling, TCP/UDP and congestion control.}
    \resumeItem{\textbf{Java \& Engineering}: JVM class loading, memory model and G1/ZGC; Synchronized, Volatile, AQS and thread-pool internals; SQL optimization, MySQL tuning and production troubleshooting.}
\resumeItemListEnd

\clearpage

% --- page 2 (中文): same engine as resume_CN.tex ---
\renewcommand{\baselinestretch}{1.14}
\titleformat{\section}
    {\sectiontitlefontsize\bfseries\color{sectiongray}}
    {}{0em}{}
\titlespacing*{\section}{0pt}{14pt}{0pt}
\renewcommand{\sectionbodyspace}{\vspace{1.5pt}{\color{applegray}\hrule height 0.8pt}\vspace{4pt}}
\renewcommand{\sectiongap}{\vspace{-3pt}}
\renewcommand{\projectgap}{\vspace{15pt}}
\renewcommand{\resumeItemListStart}{%
    \begin{itemize}[leftmargin=1.5em, labelsep=0.5em,
        label={\color{applegray}\raisebox{0.12ex}{\scriptsize\faAngleRight}},
        itemsep=6pt, parsep=0pt, topsep=1pt]}
\noindent
\begin{minipage}[t]{0.60\textwidth}
\vspace{0pt}
{\namefontsize\bfseries\color{appledark}刘龙昌}\par\vspace{4pt}
{\entrytitlefontsize\bfseries\color{textgray}软件开发工程师}\par\vspace{5pt}
{\contactfontsize\textcolor{textgray}{\faPhone}~\textcolor{textgray}{183-5757-1400}\quad \textcolor{textgray}{\faEnvelope}~\textcolor{textgray}{liulch.cn@gmail.com}}\par
\end{minipage}\hfill
\begin{minipage}[t]{0.34\textwidth}
\vspace{0pt}
\raggedright
\vspace{6pt}
{\headerLinkItem{\faGithub}{\href{https://github.com/leochame}{github.com/leochame}}}\par\vspace{3pt}
{\headerLinkItem{\faLinkedin}{\href{https://www.linkedin.com/in/longchang-liu-770211387/}{linkedin.com/in/longchang-liu}}}\par\vspace{3pt}
{\headerLinkItem{\faGlobe}{\href{https://leochame.github.io}{leochame.github.io}}}\par
\end{minipage}
\par\vspace{2pt}

\section{教育经历}
\sectionbodyspace
\begin{tabular*}{\textwidth}{l@{\extracolsep{\fill}}r}
    \textbf{\entrytitlefontsize\color{appledark}杭州电子科技大学}\hspace{1em}{\entrydetailsfontsize\color{textgray}计算机科学与技术 · 本科} & {\entrydetailsfontsize\color{applegray}2022年08月 - 2026年06月} \\
\end{tabular*}
\vspace{-14pt}
\sectiongap

\section{工作经历}
\sectionbodyspace
\begin{tabular*}{\textwidth}{l@{\extracolsep{\fill}}r}
    \textbf{\smallentrytitlefontsize\color{appledark}紫荆芯界 - Agentic Infra（Kernel-Gen Agent）} & {\entrydetailsfontsize\color{applegray}2026年07月 - 至今} \\
\end{tabular*}
{\descfontsize\color{textgray} 面向内部业务负载的 kernel 生成 agent：agent 自主完成算子生成、校验与调优，负责设计其任务契约、验证流程与评测口径，在 FlashInfer benchmark 上达到业界 SOTA 水平。}\vspace{2pt}
\resumeItemListStart
\resumeItem{\textbf{任务契约}：明确目标算子、约束与准入标准，作为 agent 每轮工作的固定输入，避免目标在长循环中漂移。}
\resumeItem{\textbf{验证流程}：每轮改动后自动与 PyTorch 参考实现数值对拍，通过后再按实测延迟与有效带宽打分，两道关卡都过才进入下一轮。}
\resumeItem{\textbf{搜索与上下文}：候选、profiling 证据与采纳结论逐轮落档，失败原因回灌为下一轮输入，只保留有效信息以控制上下文增长；配合迭代上限与性能回退拒绝，避免重复无效尝试。}
\resumeItem{\textbf{内部交付}：业务侧只需提供算子与参考实现，即可拿到通过校验的 kernel。}
\resumeItemListEnd
\projectgap

\begin{tabular*}{\textwidth}{l@{\extracolsep{\fill}}r}
    \textbf{\smallentrytitlefontsize\color{appledark}阿里巴巴 - 高德地图（实习）} & {\entrydetailsfontsize\color{applegray}2025年10月 - 2026年01月} \\
\end{tabular*}
{\descfontsize\color{textgray} 负责“订单分析助手”与“AI 打车”的大模型工程化落地，提升复杂任务处理能力并完成统一 AI 能力接入。}\vspace{2pt}
\resumeItemListStart
\resumeItem{在订单分析助手项目中，设计并落地 Deep Research 功能，提升复杂任务拆解效果与长上下文场景下的处理表现。}
\resumeItem{参与 AI 打车服务冷启动，建设统一、稳定的 AI 调用接口，保障核心链路稳定接入。}
\resumeItemListEnd
\projectgap

\begin{tabular*}{\textwidth}{l@{\extracolsep{\fill}}r}
    \textbf{\smallentrytitlefontsize\color{appledark}腾讯 - 广告 ChatBI 系统（实习）} & {\entrydetailsfontsize\color{applegray}2025年04月 - 2025年09月} \\
\end{tabular*}
{\descfontsize\color{textgray} 服务内容业务人员的自然语言取数系统：以指标中台为知识源，经多轮澄清、多路召回（reranker 重排 + 热度加权）与 DSL 逐级填充生成模板化 DSL，抑制幻觉；并以 bad case 回流 + ICL 热修复持续迭代，状态机保障宕机恢复。}\par
\sectiongap

\section{开源经历}
\sectionbodyspace
{\entrydetailsfontsize\color{textgray}\textbf{\color{appledark}LangChain4j}（12k+ Star）贡献者，Java 生态 AI 开发框架，参与社区问题修复与能力完善。}\par
\sectiongap

\section{专业技能}
\sectionbodyspace
\skillItemListStart
    \resumeItem{\textbf{AI Infra}：理解 Transformer 架构与 KV Cache；熟悉 Agent 上下文管理、Multi-Agent 协作范式、harness 设计与 eval 体系，有 GPU kernel 生成与调优实践}
    \resumeItem{\textbf{计算机基础}：扎实的操作系统与网络基础，理解内存管理（页表）、System Call、线程调度、TCP/UDP 与拥塞控制}
    \resumeItem{\textbf{Java 与工程实践}：熟悉 JVM 类加载机制、内存模型与 G1/ZGC，理解 Synchronized、Volatile、AQS 与线程池实现；有 SQL 优化、MySQL 调优与线上故障排查经验}
\resumeItemListEnd

\end{document}
